% GENERATED FILE. Edit canonical lessons, then run npm run generate:books.
% canonical-source-hash: fnv1a64:e1a8c3b0f3d61f13
% canonical-lessons: TE-C227-yuvakudu, TE-C227-peddalu, TE-C227-iruguporugu, TE-C227-satruvu, TE-C227-pellikoduku

\chapter{Young man, Elders, Neighbours, Enemy, Bridegroom}
\label{ch:te-young-man-elders-neighbours-enemy-bridegroom}
\hlchaptermodality{227}

\begin{chapteropening}
\textbf{By the end of this chapter:} \emph{I can say a young man, elders, neighbours, an enemy and a bridegroom.}

Complete the last lesson of chapter 227: I can say a young man, elders, neighbours, an enemy and a bridegroom.
\end{chapteropening}

\section[\texorpdfstring{\te{యువకుడు} (yuvakuḍu)}{యువకుడు (yuvakuḍu)}]{\te{యువకుడు} (yuvakuḍu) — a young man, a youth}

\label{lesson:TE-C227-yuvakudu}



\begin{quote}
Before the new one: say the Telugu for a form, then the Telugu for a card.
\end{quote}

\subsection*{You'll want to know: \te{యువకుడు}}
\textbf{\te{యువకుడు}} — \emph{yuvakuḍu} — \textquotedblleft{}a young man, a youth\textquotedblright{}.

\begin{grammarlens}[title={What you've built}]
One more word for this chapter.
\end{grammarlens}

\subsection*{Your turn}
\begin{itemize}
\raggedright
  \item \emph{Say it:} \emph{yuvakuḍu}
  \item \emph{Say it:} \emph{yuvakuḍu}, once more
  \item \emph{Say it:} say \emph{yuvakuḍu}
  \item \emph{Recall:} say the Telugu for a form, then the Telugu for a card, then say \emph{yuvakuḍu} again
\end{itemize}

\begin{tcolorbox}[breakable,colback=teal!4,colframe=teal!35!black,title={Before you move on}]
What does \te{యువకుడు} mean? (\textquotedblleft{}A young man, a youth\textquotedblright{}.) Say it once more.
\end{tcolorbox}
\section[\texorpdfstring{\te{పెద్దలు} (peddalu)}{పెద్దలు (peddalu)}]{\te{పెద్దలు} (peddalu) — elders}

\label{lesson:TE-C227-peddalu}



\begin{quote}
Before the new one: say the Telugu for a card, then the Telugu for a young man.
\end{quote}

\subsection*{You'll want to know: \te{పెద్దలు}}
\textbf{\te{పెద్దలు}} — \emph{peddalu} — \textquotedblleft{}elders\textquotedblright{}.

\begin{grammarlens}[title={What you've built}]
One more word for this chapter.
\end{grammarlens}

\subsection*{Your turn}
\begin{itemize}
\raggedright
  \item \emph{Say it:} \emph{peddalu}
  \item \emph{Say it:} \emph{peddalu}, once more
  \item \emph{Say it:} say \emph{peddalu}
  \item \emph{Recall:} say the Telugu for a card, then the Telugu for a young man, then say \emph{peddalu} again
\end{itemize}

\begin{tcolorbox}[breakable,colback=teal!4,colframe=teal!35!black,title={Before you move on}]
What does \te{పెద్దలు} mean? (\textquotedblleft{}Elders\textquotedblright{}.) Say it once more.
\end{tcolorbox}
\section[\texorpdfstring{\te{ఇరుగుపొరుగు} (iruguporugu)}{ఇరుగుపొరుగు (iruguporugu)}]{\te{ఇరుగుపొరుగు} (iruguporugu) — neighbours}

\label{lesson:TE-C227-iruguporugu}



\begin{quote}
Before the new one: say the Telugu for a young man, then the Telugu for elders.
\end{quote}

\subsection*{You'll want to know: \te{ఇరుగుపొరుగు}}
\textbf{\te{ఇరుగుపొరుగు}} — \emph{iruguporugu} — \textquotedblleft{}neighbours\textquotedblright{}.

\begin{grammarlens}[title={What you've built}]
One more word for this chapter.
\end{grammarlens}

\subsection*{Your turn}
\begin{itemize}
\raggedright
  \item \emph{Say it:} \emph{iruguporugu}
  \item \emph{Say it:} \emph{iruguporugu}, once more
  \item \emph{Say it:} say \emph{iruguporugu}
  \item \emph{Recall:} say the Telugu for a young man, then the Telugu for elders, then say \emph{iruguporugu} again
\end{itemize}

\begin{tcolorbox}[breakable,colback=teal!4,colframe=teal!35!black,title={Before you move on}]
What does \te{ఇరుగుపొరుగు} mean? (\textquotedblleft{}Neighbours\textquotedblright{}.) Say it once more.
\end{tcolorbox}
\section[\texorpdfstring{\te{శత్రువు} (śatruvu)}{శత్రువు (śatruvu)}]{\te{శత్రువు} (śatruvu) — an enemy}

\label{lesson:TE-C227-satruvu}



\begin{quote}
Before the new one: say the Telugu for elders, then the Telugu for neighbours.
\end{quote}

\subsection*{You'll want to know: \te{శత్రువు}}
\textbf{\te{శత్రువు}} — \emph{śatruvu} — \textquotedblleft{}an enemy\textquotedblright{}.

\begin{grammarlens}[title={What you've built}]
One more word for this chapter.
\end{grammarlens}

\subsection*{Your turn}
\begin{itemize}
\raggedright
  \item \emph{Say it:} \emph{śatruvu}
  \item \emph{Say it:} \emph{śatruvu}, once more
  \item \emph{Say it:} say \emph{śatruvu}
  \item \emph{Recall:} say the Telugu for elders, then the Telugu for neighbours, then say \emph{śatruvu} again
\end{itemize}

\begin{tcolorbox}[breakable,colback=teal!4,colframe=teal!35!black,title={Before you move on}]
What does \te{శత్రువు} mean? (\textquotedblleft{}An enemy\textquotedblright{}.) Say it once more.
\end{tcolorbox}
\section[\texorpdfstring{\te{పెళ్ళికొడుకు} (peḷḷikoḍuku)}{పెళ్ళికొడుకు (peḷḷikoḍuku)}]{\te{పెళ్ళికొడుకు} (peḷḷikoḍuku) — a bridegroom}

\label{lesson:TE-C227-pellikoduku}



\begin{quote}
Before the new one: say the Telugu for neighbours, then the Telugu for an enemy.
\end{quote}

\subsection*{You'll want to know: \te{పెళ్ళికొడుకు}}
\textbf{\te{పెళ్ళికొడుకు}} — \emph{peḷḷikoḍuku} — \textquotedblleft{}a bridegroom\textquotedblright{}.

\begin{grammarlens}[title={What you've built}]
One more word for this chapter.
\end{grammarlens}

\subsection*{Your turn}
\begin{itemize}
\raggedright
  \item \emph{Say it:} \emph{peḷḷikoḍuku}
  \item \emph{Say it:} \emph{peḷḷikoḍuku}, once more
  \item \emph{Say it:} say \emph{peḷḷikoḍuku}
  \item \emph{Recall:} say the Telugu for neighbours, then the Telugu for an enemy, then say \emph{peḷḷikoḍuku} again
\end{itemize}

\begin{tcolorbox}[breakable,colback=teal!4,colframe=teal!35!black,title={Before you move on}]
What does \te{పెళ్ళికొడుకు} mean? (\textquotedblleft{}A bridegroom\textquotedblright{}.) Say it once more.
\end{tcolorbox}
